The implementation uses one distribution adapter for every component family.
An adapter defines the unconstrained coordinates, inverse transformation, and
log density used by fitting, classification, differentiation, and inference.
Consequently, the estimated scores and decision contrasts cannot silently use
a density different from the fitted model.  Positive-definite covariance
matrices use log-Cholesky coordinates, and mixture weights use
reference-category multinomial logits.

\begin{center}
\fbox{\begin{minipage}{0.94\linewidth}
\textbf{Algorithm 1: Feasible boundary-aware covariance estimation}
\begin{enumerate}
 \item Fit the hard-classification likelihood to obtain
 $\widehat{\boldsymbol\theta}_n$.
 \item Compute the fitted labels.
 \item Compute $\widehat{\mathbf I}_{c,n}$.
 \item Compute $\widehat{\mathbf J}_n$.
 \item Compute every pairwise contrast $g_{jl}$ and its parameter gradient.
 \item Estimate $\widehat{\mathbf B}_n$ by \eqref{eq:boundary-estimator}.
 \item Form $\widehat{\mathbf A}_n=\widehat{\mathbf I}_{c,n}-
 \widehat{\mathbf B}_n$ and stabilize it if necessary, recording the diagnostic.
 \item Report $\widehat{\mathbf V}_n=\widehat{\mathbf A}_n^{-1}
 \widehat{\mathbf J}_n\widehat{\mathbf A}_n^{-\mathsf T}$ and standard errors
 from $\widehat{\mathbf V}_n/n$.
\end{enumerate}
\end{minipage}}
\end{center}

For the multivariate skew-normal component, we use
\begin{equation}
 f(\mathbf{x};\boldsymbol{\xi},\boldsymbol{\Omega},\boldsymbol{\alpha})
 =2\phi_p(\mathbf{x};\boldsymbol{\xi},\boldsymbol{\Omega})
 \Phi\!\left\{
 \boldsymbol{\alpha}^{\mathsf T}\boldsymbol{\Omega}^{-1/2}
 (\mathbf{x}-\boldsymbol{\xi})
 \right\},
 \label{eq:skew-normal-density}
\end{equation}
where $\boldsymbol{\Omega}^{-1/2}$ is the symmetric inverse square root.  The
same expression is used by the fitting library and the inference code.

Hard-EM fitting explicitly requests the hard mode.  We use both K-means and
Gaussian-mixture initializations.  When component families are distinct, all
family-to-component assignments are fitted and the result with the largest
empirical classification likelihood is retained.  The objective gap, selected
initialization, component weights, tail parameter, covariance eigenvalues, and
skewness norm are retained as basin-stability diagnostics.  The result is then
mapped to the manuscript's canonical family order.  Repeated families are
ordered by a prespecified fitted location functional.  Neither operation uses
the generating parameters or simulated labels.

The reference implementation evaluates derivatives by centered finite
differences in the common unconstrained coordinate system.  Evaluations are
vectorized over observations, so a full contrast-gradient matrix requires
two score evaluations per parameter coordinate rather than two evaluations
per observation and coordinate.  A deliberately slow observation-wise
implementation is retained as an independent numerical check.

The contrast $g_{jl}$ is a dimensionless log posterior odds.  Our prespecified
pilot bandwidth therefore uses its natural scale,
\begin{equation}
 h_{jl}=1.06n^{-1/5}.
 \label{eq:pilot-bandwidth}
\end{equation}
The global empirical spread of $g_{jl}$ is reported but does not multiply the
bandwidth: in separated mixtures it is dominated by observations far from the
boundary and produces severe oversmoothing. This is a pilot rule, not an
optimality claim. Every analysis reports $h_{jl}$, the robust global spread,
the number of observations for which $(j,l)$ is active, and the number falling
inside the kernel support.
Pairs with insufficient effective observations are flagged rather than
silently assigned a wider bandwidth.  The simulation study examines
multiplier sensitivity around \eqref{eq:pilot-bandwidth}.

We report the raw corrected matrix and its nonpositive-curvature failure rate.
Primary finite-sample inference uses the eigenvalue-constrained version in
\eqref{eq:boundary-stabilization}, with $\varepsilon_n=n^{-1/4}$.  The
associated shrinkage factor is reported for every replication; the correction
is asymptotically unmodified by Proposition~\ref{prop:stabilization-inactive}.

Population CML targets used in simulation are approximated from several
independent large reference samples.  We retain every resulting target vector
and report the between-reference Monte Carlo standard error of each coordinate.
This prevents numerical target uncertainty from being treated as zero.  Target
fits and ordinary simulation fits use the same multistart hard-EM,
classification-objective selection, and canonical-ordering rules.

The nonparametric bootstrap used to assess local inferential accuracy is
initialized at the original estimate and constrained to a local trust region.
Replicates reaching that region's boundary are reported as failures.  Fresh
global bootstrap fits that reselect the family assignment answer a different
question---stability under model and basin selection---and are reported
separately rather than pooled into a Gaussian standard error.
